\section{Related Work and Deep Learning Literature Review}
\label{sec:related}

\subsection{From statistical models to feature-based learning}
Ahmed and Cook~\cite{ahmed1979} first applied Box--Jenkins ARIMA to freeway data,
and Williams and Hoel~\cite{williams2003} justified seasonal ARIMA by showing that
weekly-differenced traffic is stationary. These models are univariate and linear,
and they degrade quickly with horizon: on METR-LA, ARIMA's MAE rises from 3.99 to
6.90\,mph between 15 and 60\,min~\cite{li2018dcrnn}. Non-parametric
regression~\cite{smith1997}, random forests~\cite{breiman2001} and gradient-boosted
trees~\cite{friedman2001,chen2016xgboost,ke2017lightgbm} capture non-linear
interactions and remain strong per-segment forecasters. Their accuracy, however,
depends on hand-built features, and a model fitted to one segment cannot see what is
happening on the links around it.

\subsection{Deep sequence models}
Stacked autoencoders~\cite{lv2015} and LSTMs~\cite{hochreiter1997,ma2015lstm}
learn temporal representations directly from raw series, and encoder--decoder
LSTMs~\cite{sutskever2014} forecast a whole horizon at once (4.37\,mph at 60\,min on
METR-LA~\cite{li2018dcrnn}). Gated convolutions~\cite{dauphin2017glu} and dilated
causal convolutions~\cite{oord2016wavenet} reach the same receptive field in
parallel. All of these treat sensors as independent sequences or as one flat vector,
so road topology is either ignored or must be learned from scratch.

\subsection{Graph neural networks for traffic}
Graph convolutions generalise convolution to irregular domains, either
spectrally~\cite{defferrard2016} or with first-order neighbour
aggregation~\cite{kipf2017}. GraphSAGE~\cite{hamilton2017sage} replaces the fixed
normalisation with learnable aggregators such as the neighbour mean; a survey is
given in~\cite{wu2021gnnsurvey}. DCRNN~\cite{li2018dcrnn} models traffic as
bidirectional diffusion on the directed distance graph inside a GRU sequence-to-sequence
model. STGCN~\cite{yu2018stgcn} interleaves gated temporal convolutions with
Chebyshev graph convolutions in a ``sandwich'' block, and Graph
WaveNet~\cite{wu2019gwnet} adds a self-adaptive adjacency learned from node
embeddings. MTGNN~\cite{wu2020mtgnn} and AGCRN~\cite{bai2020agcrn} learn the graph
entirely. STID~\cite{shao2022stid} shows that simple spatial and temporal identity
embeddings recover much of the gain, which motivates the node embeddings in our model.
Jiang and Luo~\cite{jiang2022survey} survey the area.

\subsection{Attention on graphs}
Attention~\cite{bahdanau2015,vaswani2017} computes data-dependent weights over a
set of inputs. Graph Attention Networks (GAT)~\cite{velickovic2018gat} apply it to
a node's neighbourhood. The coefficient $\alpha_{ij}$ that node $i$ assigns to
neighbour $j$ is a softmax over learned scores
$\operatorname{LeakyReLU}(\vect{a}^\top[W\vect{h}_i\,\|\,W\vect{h}_j])$, so each
neighbour's influence depends on the current states rather than on fixed edge weights.
GATv2~\cite{brody2022gatv2} shows that this original form is ``static'': the ranking
of neighbours is the same for every query node. It proposes a more expressive
variant, which we discuss as future work. In traffic, GaAN~\cite{zhang2018gaan}
gates attention heads, ASTGCN~\cite{guo2019astgcn} learns spatial and temporal
attention matrices shared across an input window, and ST-GRAT~\cite{park2020stgrat}
and GMAN~\cite{zheng2020gman} build encoder--decoders from spatio-temporal attention.
Attention suits traffic because the relevance of a neighbour changes with
conditions: a downstream queue matters only when one is forming.

\subsection{Hybrid statistical / machine learning / deep models}
Combining a strong structured model with a neural network has a long history in
forecasting. Zhang~\cite{zhang2003hybrid} fits ARIMA to the linear component and
a neural network to its \emph{residuals}. The M4-winning ES-RNN~\cite{smyl2020esrnn}
couples exponential smoothing with an LSTM. Stacked
generalisation~\cite{wolpert1992stacking} trains a second-level learner on the
first-level model's \emph{out-of-sample} predictions, because in-sample predictions
are optimistically biased and teach the meta-learner the wrong residual. For
dependent data the folds must be contiguous blocks~\cite{bergmeir2012cv}, and
leakage of future information is a well-documented failure
mode~\cite{kaufman2012leakage}. Residual learning~\cite{he2016resnet} shows that
learning a correction to a good reference is easier than learning the full mapping.
Our hybrid combines these ideas. A per-segment gradient-boosting model supplies the
reference, and a graph attention network learns the network-level residual. The
blocked cross-fitting follows~\cite{bergmeir2012cv,wolpert1992stacking}.

\subsection{Missing data and regularisation}
Treating an absent reading as a value biases a network. GRU-D~\cite{che2018grud}
shows that explicit missingness masks carry information and improve accuracy, and we
adopt such a mask channel. Dropout~\cite{srivastava2014dropout}, decoupled weight
decay~\cite{loshchilov2019adamw}, layer normalisation~\cite{ba2016layernorm} and
early stopping~\cite{prechelt1998} are the standard regularisers for models of
our size.

\subsection{Signal control and simulation}
Forecasts become valuable when they change a control decision. Classical ramp
metering such as ALINEA~\cite{papageorgiou1991alinea} reacts to \emph{measured}
occupancy. Reinforcement-learning signal controllers (e.g.\
IntelliLight~\cite{wei2018intellilight}) are trained and evaluated in microscopic
simulators such as SUMO~\cite{krajzewicz2012sumo} and CityFlow~\cite{zhang2019cityflow}.
METR-LA provides speeds only, with no turning movements or signal timings. We
therefore restrict ourselves to an \emph{offline, simulated} plan-switching
experiment that measures whether a forecast lets a controller act before congestion
arrives. We make no claim of live control.

\textbf{Positioning.} Most published GNN comparisons on METR-LA pit
sophisticated graph models against feature-poor baselines and do not test
significance. We instead (i) build the strongest per-segment baseline we can,
(ii) measure the specific information it lacks, (iii) fill exactly that gap with
graph attention over the fixed road graph, and (iv) quantify every design choice by
ablation and Diebold--Mariano tests.
